% 第 7 章。数字：checks/check_ch07.py；引文：data/snapshots/greek_facts.csv（逐字核对见
% data/fetch/verify_web_quotes.py）；埃及的价格：kelly_egypt_prices；图：figsrc/fig_ch07.py（greek_trust 与 greek_ceiling 是教学模型）。
\chapter{希腊与罗马：铸币、城邦财政与帝国的贬值}\label{ch:greekroman}

\begin{chapterguide}
\guidequestion{罗马的银币为什么越铸越薄？古代经济算不算市场经济？}
\begin{guidegoals}
\item 说明铸币的出现怎样改变了交易与财政（吕底亚铸币、雅典劳里昂银矿）；
\item 用罗马银币含银量的下降与三世纪的价格上涨说明铸币贬值的机制；
\item 分析戴克里先限价令（301 年）为什么失败；
\item 概括 Finley 的“原始主义”与 Temin 的“市场经济”两种解读，以及双方的证据。
\end{guidegoals}
\guideprereq{\chref{ch:mesopotamia}\enspace\nameref{ch:mesopotamia}；\chref{ch:prices}\enspace\nameref{ch:prices}}
\end{chapterguide}

希罗多德（Herodotus）说，吕底亚人（Lydians）“是我们所知道的最早铸造和使用金币、银币的人，也是最早的零售商”\cite{herodotus}。上一章的白银要称，大麦要量；金属块上打了印记以后，钱可以数了。

到了罗马帝国治下的埃及，纸草文书记下了许多小麦的价格。粮食按阿塔巴（artaba）量，这是埃及的谷物容量单位；价钱按德拉克马（drachma）算，这是埃及记账用的货币单位。按 Kelly（2021）整理的数据，公元 270 年，一份庄园账目里的小麦每阿塔巴 24 德拉克马。276 年，官府把实物税折成钱，按每阿塔巴 200 德拉克马算。270 年到 276 年，小麦的价格涨到原来的 8.3 倍。约 340 年，奥克西林库斯（Oxyrhynchus）的商人申报了七个价格，中位数是 45 塔兰特（talent）。塔兰特是比德拉克马大得多的记账单位，这个价钱合 27 万德拉克马\cite{kelly2021}。

钱本身也变了。罗马的银币叫第纳尔（denarius），尼禄（Nero）在 64 年把它的含银定为 80\%。到 260 年代末，帝国的“银”币含银已不到 5\%\cite{haklai2011}。

物价却没有亦步亦趋。Haklai-Rotenberg（2011）说，零散的价格证据“并没有明确表明通胀紧随银币的贬值而来”。274 年，奥勒良（Aurelian）发行了新币。钱币的含银从大约 2\% 或 3\% 提高到将近 5\%，此前相对平稳的物价却“突然上涨了十倍”\cite{haklai2011}。

301 年，戴克里先（Diocletian，284—305 年在位）颁布限价令\cite{kent1920}。拉克唐修（Lactantius）在《论迫害者之死》里说，此后为了最微不足道的东西流了许多血，人们不敢把任何东西拿出来卖，短缺比以前更加严重\cite{lactantius}。

银币为什么越铸越薄，银少了物价为什么没有马上涨，274 年钱变好了物价为什么反而暴涨，限价令又为什么管不住物价？

\section{当时的世界}\label{sec:greek-world}

开场的几件事，分别发生在吕底亚与雅典、罗马与罗马治下的埃及，以及戴克里先的帝国。关于它们，今天凭的是哪些史料，这些史料又能说明什么？

\subsubsection{吕底亚与雅典}
希罗多德没有说吕底亚人从什么时候开始铸币。雅典的货币单位，《雅典政制》（英译本题为亚里士多德所作）归于梭伦（Solon）的改革。梭伦把原来以 70 德拉克马为标准的弥那（mina），提高到足足 100 德拉克马。他还规定了与钱币相应的重量单位，“63 弥那合一塔兰特，多出的三弥那分摊到斯塔特与其他币值上”\cite{athpol}。本章用的是货币单位，不用这套重量单位，1 弥那 = 100 德拉克马。1 塔兰特 = 60 弥那 = 6000 德拉克马。Kelly 也说，每塔兰特等于 6000 银德拉克马。更小的单位是奥波尔（obol），1 德拉克马 = 6 奥波尔，这是 Kelly 对埃及 190 年以前的折算\cite{kelly2021}，色诺芬的数字也合这个比例（例~\ref{ex:greek-xenophon}）。

雅典附近有银矿。《雅典政制》说，在尼科德穆斯（Nicodemus）任执政官的那一年，“马罗尼亚（Maroneia）的矿被发现了，国家从开采中得到一百塔兰特的利润”。三年以后，在许普西基德斯（Hypsichides）任执政官那一年，由于薛西斯（Xerxes）的大军逼近，被放逐的人都被召回了\cite{athpol}。希罗多德说的是劳里昂（Laurion，他的英译本写作 Laureion）的银矿\cite{herodotus}。这些文献用执政官的名字与薛西斯的进军来纪年，没有给出公元前的年份，这里也不折算。这笔钱怎样花掉，是\secref{sec:greek-laurion}的故事。

色诺芬（Xenophon，公元前 431—前 354 年）的《论收入》，是给雅典出的增加公共收入的主意，一大半讲的是银矿\cite{xenrev}。

\subsubsection{罗马的银币与埃及的纸草}
肯特（R. G. Kent）说，在共和国晚期与帝国早期，标准的钱币是银第纳尔\cite{kent1920}。埃及在罗马治下另有一套钱。Howgego、Butcher、Ponting 与 Heuchert（2010）说，“罗马埃及实行一种封闭的货币制度”。那里流通的是含银不高的四德拉克马（tetradrachm），每枚与一第纳尔等值\cite{howgego2010}。Kelly 说，这种在亚历山大里亚铸造的四德拉克马，一枚值四德拉克马\cite{kelly2021}。埃及的物价因此用德拉克马计。

为什么用埃及？因为别处几乎没有价格。Howgego 等说，Rathbone 对埃及价格的研究，是整个帝国唯一可用的价格序列\cite{howgego2010}。本章的价格来自 Kelly（2021）的附录 2。Kelly 从纸草文书里辑出各种商品的价格，一份文书一行\cite{kelly2021}。本章用其中五张表，三张在 274 年以前。私人小麦价格有 40 份文书，国家小麦价格有 31 份，大麦价格有 25 份。另两张在 274 年以后，小麦有 23 份，大麦有 11 份。五张表共 130 份文书。照 Kelly 的做法，一份文书里有几个价格时取中位数，年代不确切的文书假定在年代范围的中点。奥波尔在 190 年以后按 7 个、以前按 6 个折一德拉克马\cite{kelly2021}。\secref{sec:greek-test} 末尾的“深入一步”还会回到这些处理。

\subsubsection{戴克里先的帝国}
这道限价令，肯特 1920 年在《宾夕法尼亚大学法律评论》上介绍过。他说，今天知道戴克里先是一个出身低微的达尔马提亚士兵。284 年，戴克里先被自己的士兵拥立为奥古斯都，即罗马皇帝。286 年，他任命了一位共治的奥古斯都，293 年又任命了两位恺撒作为助手\cite{kent1920}。

从此，帝国要养的人多了。拉克唐修说，军队成倍增加，四位君主每一位都力图维持一支比过去任何一位独掌帝国的皇帝都大得多的军队；“领工钱的人开始比纳税的人还多”\cite{lactantius}。肯特提醒，拉克唐修“是一个基督徒，对迫害他的教友的政府心怀怨恨”，很难说心态公正。但肯特也同意，如今要维持的不只是一个皇宫，而是四个，每个都有自己的官员和军队\cite{kent1920}。301 年，戴克里先颁布限价令。305 年，即限价令颁布四年以后，他退位了\cite{kent1920}。

\section{铸币：印记省掉了称量}\label{sec:greek-coin}

吕底亚人打在金属上的印记，到底改变了什么？一枚钱币的价值，来自它所含的金属，还是来自约定？亚里士多德的书里，两种回答都有。

\subsubsection{数而不称}
亚里士多德（Aristotle）在《政治学》里讲钱的起源。人们发明了一种东西来彼此交换，大家都愿意给，也都愿意收；它本身确有价值，还另有便于携带的好处。这种东西“起初只按重量或大小计值，但随着时间推移，人们给它打上一个印记，省去称量的麻烦，印记表明它的价值”\cite{politics}。这与\chref{ch:mesopotamia}的世界只差一步。那里的白银按重量计价，每一笔付款都要称；有了印记，钱就可以数。印记是发行者的担保，它省掉的是称量与验成色的成本，这是交换成本的一部分。

可是“印记表明它的价值”，表明的是谁的价值？如果印记只是证明了金属的重量与成色，钱的价值仍然是金属的价值。一旦大家习惯了数而不称，发行者就有了一个机会，少放一点银子，印记照旧。这一章的大半，讲的就是发行者怎样用掉这个机会。

\subsubsection{约定还是金属}
亚里士多德在《尼各马可伦理学》里给出了另一种回答。他说，钱“之所以是它，不是出于自然，而是出于习俗或法律（希腊文 νόμος），改变它的价值，或者让它完全无用，都在我们”。他又说，钱是一种保证，让我们将来可以交换。它当然也会贬值，因为它的购买力并不总是一样，但它比它所代表的商品更持久\cite{ethics}。他还用一个例子说明钱是计价的尺度：“为一所房子付五张床，还是付五张床的价钱，是没有分别的”\cite{ethics}。这正是\chref{ch:mesopotamia}的计价单位。

两段话的侧重不同。《政治学》的这一段讲钱的来历，被选来交换的东西本身确有价值，印记只是省掉了称量。《伦理学》的这一段讲钱的性质，它出于约定，可以改变。《政治学》在下文也转述了后一种看法，说另一些人“认为它毫无价值，因为它天然就没有价值，只是凭约定才有了价值；所以用它的人一旦改变主意，它就一文不值”\cite{politics}。后来的人把前一种说法叫作金属论，后一种叫作名目论。\secref{sec:greek-prices} 的模型把两者放进同一个式子。一枚钱币值多少，介于它所含的金属与它的法定价值之间；落在哪里，取决于人们对它的信任。

\subsubsection{出了边境，只是白银}
色诺芬注意到，同一枚钱在边境内外的命运不同。他说，在大多数城邦，商人卖了货，还得装上别的货物运回去，因为“现行的钱币出了边境就不流通”。在雅典却不必，“他只要把白银运出去就行，他找不到比这更好的出口货物，因为不论在哪里卖，他都可以指望在本钱上赚一大笔”\cite{xenrev}。边境以内，钱按约定数着用；边境以外，钱只是一块银子，按重量与成色作价。雅典的钱含银足，出了边境照样值钱，所以白银本身就成了雅典的出口货。有的钱只在本国通行，有的钱（或它所含的金属）到哪里都有人要，这是\chref{ch:hierarchy}要讲的货币等级的一个古老的例子。

\begin{intuition}[印记是一种担保]
一块没有印记的银子，每次易手都要称、要验；一枚有印记的钱，收钱的人只看印记。印记越可靠，收钱的人越省事，钱也就越好用。可是收钱的人既然不再称、不再验，他就不知道钱里实际有多少银子；他相信的是印记背后的发行者。这份信任有两面：它让钱比一块银子更好用，也让发行者可以少放银子而不被马上发现。信任用得越多，剩下的就越少。
\end{intuition}

\section{劳里昂的白银：一笔意外之财}\label{sec:greek-laurion}

雅典的白银出了边境也值钱，城邦还从银矿得到过一笔意外之财。这样一笔钱，该分掉，还是用掉？靠银矿，能不能长期养活一个国家？

\subsubsection{分掉，还是造船}
希罗多德说，雅典人的国库里有一大笔从劳里昂的银矿收进来的钱，他们打算把它分掉，每人领十德拉克马。地米斯托克利（Themistocles）说服雅典人放弃分钱，用这笔钱造两百艘船去打仗，他说的是同埃伊纳人的战争\cite{herodotus}。希罗多德认为，这场战争“在当时确实拯救了希腊，因为它迫使雅典人成为一个海上强国”\cite{herodotus}。

《雅典政制》的说法更细。有人建议人民把这笔钱分掉，地米斯托克利阻止了。他不肯说打算把钱花在什么上，只叫人民把钱借给雅典最富有的一百个人，每人一塔兰特。条件是：“如果钱的用法让人民满意，这笔开支就由国家承担；否则国家就从借钱的人那里把钱收回”\cite{athpol}。拿了钱的人每人造了一艘三列桨战船，一共一百艘，雅典人就是用这些船在萨拉米斯（Salamis）同蛮族作战的\cite{athpol}。

这是一份很巧妙的合同。国家把钱交给最富的一百个人去办一件具体的事。办得好，钱算国家出的；办不好，钱要还。富人们不能拿了钱不办事，因为办不好就要自己掏钱；国家也不必先说清楚钱要用在哪里。用今天的话说，地米斯托克利把一笔资源收入变成了公共投资，并用一个依结果而定的条款约束了执行的人。一笔资源收入该分给国民，还是用来投资，今天的产油国仍在争论（\chref{ch:oil}）。

\begin{example}{一百塔兰特是多少钱}\label{ex:greek-laurion}
按《雅典政制》，马罗尼亚的矿给国家带来一百塔兰特的利润，造了一百艘战船；按希罗多德，原来的打算是每人分十德拉克马，后来造了两百艘船。希罗多德在讲另一件事时提到过“三万雅典人”\cite{herodotus}。（1）一百塔兰特是多少德拉克马？按每人十德拉克马，够分给多少人？（2）如果按三万人每人十德拉克马来分，要多少塔兰特？（3）按《雅典政制》，一艘战船花多少？照这个价钱，希罗多德的两百艘船要多少？
\solution （1）$100\times6000=60$ 万德拉克马，按每人 10 德拉克马，够分给 6 万人，是“三万雅典人”的 2 倍。（2）$30000\times10=30$ 万德拉克马，合 50 塔兰特，只是马罗尼亚利润的一半。（3）$100/100=1$ 塔兰特一艘；两百艘就要 200 塔兰特，是马罗尼亚利润的 2 倍。两种记载对不上：或者说的不是同一笔钱（两者都说钱来自银矿，希罗多德只说“一大笔钱”），或者分钱的人数、船的价钱另有不同。“三万”是希罗多德在别处用的一个整数，不是户籍统计。这些数字能说明的是数量级：这笔银矿收入，够造一支一百艘到两百艘船的舰队，也够给每个公民分十德拉克马。
\end{example}

\subsubsection{色诺芬的国营奴隶}
色诺芬想让银矿长期养活雅典。他的论证从需求开始。发现的矿石越多，炼出的白银越多，愿意干这一行的人也越多。别的东西多了就卖不掉，白银不一样：“白银，从来没有人多到不得不喊‘够了’的地步”。用不完的人把它埋进地里藏起来，同实际用它一样快活\cite{xenrev}。

他举了几个私人的例子。尼刻拉图斯之子尼西阿斯（Nicias）在银矿里拥有一千人，把他们租给索西亚斯（Sosias）。索西亚斯为这一千个奴隶中的每一个，每天付给他净一个奥波尔，不扣任何费用。希波尼库斯（Hipponicus）有六百个奴隶，照同样的办法出租，每天净得一弥那；菲勒摩尼德斯（Philemonides）有三百个，每天带来半弥那\cite{xenrev}。

色诺芬的建议是由国家照样去做，置办“一批公共奴隶，数目比如说是每个雅典公民三个”，租给开矿的人。他设想起初先组建一千二百个奴隶的核心，单靠利润，五六年之内把人数增加到至少六千。“这六千人，假设每个奴隶每天净挣一个奥波尔，我们一年就有六十塔兰特的收入”；到了一万这个整数，年收入就将达到一百塔兰特\cite{xenrev}。

\begin{example}{色诺芬的算术}\label{ex:greek-xenophon}
（1）尼西阿斯的一千个奴隶，每人每天一个奥波尔，他一天收多少德拉克马？一年（按 360 天）多少塔兰特？（2）希波尼库斯的六百个奴隶每天带来一弥那，菲勒摩尼德斯的三百个带来半弥那。这两个数与“一德拉克马合六奥波尔、一弥那合一百德拉克马”相合吗？（3）色诺芬说六千个奴隶一年带来六十塔兰特。他的一年是多少天？按 365 天算是多少？（4）如果真的“每个公民三个”，按希罗多德的“三万雅典人”，要多少奴隶？一年收入多少？
\solution （1）一天 1000 奥波尔，即 $1000/6=166.7$ 德拉克马；一年 $166.7\times360/6000=10$ 塔兰特。（2）600 奥波尔 $=100$ 德拉克马 $=1$ 弥那，300 奥波尔 $=50$ 德拉克马 $=0.5$ 弥那，正好相合；色诺芬的数字也印证了这两个换算。（3）60 塔兰特合 36 万德拉克马，即 216 万奥波尔，除以 6000 个奴隶，是 360 天；按 365 天算是 60.8 塔兰特。色诺芬用的是一个取整的年。同样，一万个奴隶一年 $10000\times360/6/6000=100$ 塔兰特，与他说的一样。（4）9 万个奴隶，一年 900 塔兰特，是马罗尼亚那笔意外之财的 9 倍。这只是按他的比例推算；他自己的计划只说到“一万这个整数”。
\end{example}

这个计划有两个前提。一是租金。每个奴隶每天净一个奥波尔，要求矿上永远需要人手，矿石永远采不完。二是奴隶本身。在色诺芬的账上，他们是按天出租的资产，每人每天净挣的一个奥波尔全归主人；可他们是人，一有机会就会逃走。英译者 Dakyns 在注里说，色诺芬指的是公元前 413 年狄凯里亚（Decelea）被设防以后的事。据修昔底德（Thucydides），那时“两万多个奴隶逃走了，其中许多是手艺人”\cite{xenrev}。战事一来，这笔收入连同本金可以一起消失。这个计划后来是否实行过，本章的史料没有说。

\subsubsection{白银会不会贬值}
色诺芬说，矿多了，白银也不会不值钱，既不用担心矿石会采完，也不用担心白银会贬值。他知道黄金会：“黄金一旦大量涌入，贬值的是黄金本身，同时还会使白银的价值上升”\cite{xenrev}。

他的论证有一个漏洞。他说白银不会贬值，理由是人们对白银的需求没有止境，用不了就埋起来。可是对黄金也可以这样说，他却承认黄金会贬值。一种金属多了而价值不降，需要人们愿意按原来的价钱无限地收下它。只要这份需求有止境，供给大增，它以别的东西计的价值就要下跌。美洲的白银大量流入欧洲与中国以后，银价与物价怎样变化，是\chref{ch:silver}的主题（式~\eqref{eq:silver-qty}）。色诺芬也有说对的地方。白银的需求里，有很大一部分是窖藏，而窖藏多少，会随人们对将来的看法变化。人们对将来的看法，会不会也左右一枚钱值多少？

\section{罗马的银币怎样变薄}\label{sec:greek-denarius}

罗马的第纳尔起初九成是银。到三世纪后期，所谓银币，含银已不足一成。含银是怎样一步步降下来的？皇帝为什么要这样做，从中能得到多少？

\subsubsection{从九成到不足一成}
表~\ref{tab:greek-silver} 把本章用到的化验与研究排在一起。读这张表要注意两点。第一，Cope（1974）用的是罗马的重量单位。一罗马磅（libra）合 12 罗马盎司（uncia）\cite{kent1920}，一盎司合 24 斯克鲁普尔（scrupulum）\cite{cope1974}，一磅就是 288 斯克鲁普尔。Cope 的标准写成“每磅银加 $x$ 盎司合金”时，含银是 $12/(12+x)$；写成“每磅含银 $s$ 斯克鲁普尔”时，含银是 $s/288$。第二，各家的方法与样本不同，数字不完全一致。塞普蒂米乌斯·塞维鲁（Septimius Severus）的大贬值，Cope 的标准折合 44.4\%，Haklai-Rotenberg 说是 50\% 左右。伽利埃努斯（Gallienus）时，Cope 的标准是 1\%—2\%。Haklai-Rotenberg 则说，260 年代末不到 5\%。表中的安东尼尼安努斯，是卡拉卡拉推出的一种比第纳尔重的新银币。

\begin{table}[htbp]\centering\small
\caption{罗马银币与埃及四德拉克马的含银}\label{tab:greek-silver}
\begin{tabularx}{\textwidth}{@{}llYl@{}}
\toprule
年代 & 钱币 & 含银 & 依据 \\
\midrule
64 年以前 & 第纳尔 & 约 90\% & Haklai-Rotenberg \\
64 年（尼禄） & 第纳尔 & 80\%（银四铜一）；此后新旧两种标准并用 & Haklai-Rotenberg \\
100 年（图拉真） & 第纳尔 & 80\%，尼禄的标准“完全确立” & Haklai-Rotenberg \\
尼禄以后 & 第纳尔 & 92.3\% 逐步降到 70.6\%（每磅银加 1 盎司合金，逐次多加 1 盎司，直到 5 盎司） & Cope \\
到 155 年 & 第纳尔 & 80\% & Howgego 等 \\
155 年起 & 第纳尔 & 70\% & Howgego 等 \\
180 年起 & 第纳尔 & 67\% & Howgego 等 \\
193 年（塞维鲁） & 第纳尔 & 44.4\%（每磅银加 15 盎司合金）；约 50\% & Cope；Haklai-Rotenberg \\
211—217 年（卡拉卡拉） & 安东尼尼安努斯 & 重约 1.5 第纳尔，作价大概 2 第纳尔；合金与同时的第纳尔相同（44.4\%） & Haklai-Rotenberg；Cope \\
伽利埃努斯独掌帝位时 & 安东尼尼安努斯 & 1.04\%，后来 2.08\%（每磅 3、后来 6 斯克鲁普尔）；260 年代末不到 5\% & Cope；Haklai-Rotenberg \\
274 年（奥勒良） & 新币 & 约 2\%—3\% $\to$ 将近 5\% & Haklai-Rotenberg \\
到约 294 年 & 帝国的“银”币 & 以贱金属为主 & Cope \\
\midrule
到 167 年 & 埃及四德拉克马 & 18\%（每枚含银 2.3 克） & Howgego 等 \\
167/8—168/9 年 & 埃及四德拉克马 & 5\%（0.6 克）；最早的证据在 164/5 年 & Howgego 等 \\
169/70 年起 & 埃及四德拉克马 & 10\%（1.3 克） & Howgego 等 \\
185/6 年起 & 埃及四德拉克马 & 降到 5\%（0.6 克） & Howgego 等 \\
\bottomrule
\end{tabularx}
\datasource{Haklai-Rotenberg（2011），第 7—8、16 页；Cope（1974），第 103、106—107、124—125 页（表 III；表 X：24 斯克鲁普尔 = 1 盎司）；Howgego、Butcher、Ponting 与 Heuchert（2010），第 8、11 页（表 1—3）。Cope 的标准按 $12/(12+x)$ 与 $s/288$ 换算成百分数，计算见 \texttt{checks/check\_ch07.py}。}
\end{table}

从头到尾看，第纳尔的含银从 80\% 降到 2\%，同样重的一枚钱，银子少了 97.5\%。埃及的四德拉克马降得早。167 年以前，每枚含银 2.3 克。到 185/6 年以后，只有 0.6 克。含银少了 74\%，只剩原来的 $1/3.83$。同一时期，第纳尔每枚的含银从 2.4 克降到 1.9 克，少了 21\%\cite{howgego2010}。

\subsubsection{为什么要贬值}
Howgego 等说，白银供应减少，加上军费的需要，是铸币贬值的典型背景\cite{howgego2010}。拉克唐修抱怨领工钱的人比纳税的人还多，说的是同一件事的财政一面。开支涨得比税收快，缺口总要有东西来填。埃及四德拉克马的大贬值在 164/5 年开始，那时瘟疫还没有到来。到 167/8 年，贬值已全面展开。Howgego 等认为，它的时机与幅度“很可能是帕提亚战争的财政压力逼出来的，瘟疫在多大程度上加重了它，无从知道”\cite{howgego2010}。

贬值可以做得不动声色。Haklai-Rotenberg 引 Butcher 与 Ponting 的话说，“过去的铸币当局是掩饰其钱币含银下降的行家”\cite{haklai2011}。也可以换一种钱。卡拉卡拉（Caracalla，211—217 年在位）推出一种新币，学术文献常称为安东尼尼安努斯（antoninianus）。它的重量约为第纳尔的 1.5 倍，法定价值大概是两个第纳尔\cite{haklai2011}。

史家卡西乌斯·狄奥（Cassius Dio）说得更难听。卡拉卡拉拿去收买易北河口一带部落的“金子当然是真的，而他给罗马人的银币和金币却是成色不足的：前者用铅镀银制成，后者用铜镀金制成”\cite{dio78}。狄奥的说法要与化验对照。按表~\ref{tab:greek-silver}，塞维鲁的标准下第纳尔含银四成多，降到百分之几要到 260 年代。铅芯镀银的钱，在本章引用的化验里没有出现。Cope 的化验说，214 年以后，同时的第纳尔与安东尼尼安努斯用的是同一种合金成色\cite{cope1974}。按这个结果，卡拉卡拉的安东尼尼安努斯含银仍是塞维鲁标准的 44.4\%。狄奥也许夸大了，也许说的是别的钱，例如镀银的伪币，本章的史料不能判断。能确定的是，那时已经有人觉得钱在变坏。

\begin{example}{皇帝从贬值里得到多少}\label{ex:greek-seign}
\chref{ch:medieval}的式~\eqref{eq:medieval-seigniorage} 定义了毛铸币税率 $1-MP/ME$：$ME$ 是一磅金属能铸出的钱的面值，$MP$ 是铸币厂为一磅金属付出的钱。设钱币的重量不变，铸币厂收银子时仍按旧钱的含银付价（$MP$ 等于贬值前的 $ME$）。（1）塞维鲁把标准从每磅银加 5 盎司合金改为加 15 盎司。新钱的含银是旧钱的几分之几？$ME$ 提高多少？毛铸币税率是多少？（2）卡拉卡拉的安东尼尼安努斯重 1.5 第纳尔，作价 2 第纳尔；按 Cope 的化验，它与同时的第纳尔用同一种合金\cite{cope1974}。每一第纳尔面值所含的银是原来的几分之几？毛铸币税率是多少？（3）两次合起来，每一第纳尔面值的含银是 193 年以前的几分之几？
\solution （1）含银从 $12/17=70.6\%$ 降到 $12/27=44.4\%$，新钱是旧钱的 $44.4/70.6=0.63$；同样一磅银能铸出的钱是原来的 $1/0.63=1.59$ 倍，$ME$ 提高 59\%。按式~\eqref{eq:medieval-seigniorage}，毛铸币税率是 $1-0.63=37\%$：铸币厂每收进一磅银，铸出的钱里有 37\% 归皇帝。（2）每一第纳尔面值的含银是 $1.5/2=0.75$，毛铸币税率 25\%。作价 2 第纳尔是多数学者的看法；另有一种看法“认为这种新币的法定价值是 1.5 第纳尔”\cite{haklai2011}，那样每一第纳尔面值的含银不变，铸币税是零。（3）按作价 2 第纳尔，$0.63\times0.75=0.47$，不到原来的一半。这里的前提是交银子的人仍按旧价交（银子若来自国库，例如把收上来的旧钱重铸后发给军队，皇帝得到的同样是少放的这部分银子）；如果人们看穿了，要求按新钱的含银付价，$MP$ 跟着 $ME$ 一起提高，铸币税就没有了。皇帝能得到多少，取决于人们多久才看穿，这正是下一节的问题。
\end{example}

铸币贬值与今天的纸币超发，做的是同一件事，都是国家用发行货币来填补开支的缺口。两者的差别在于约束。

\begin{compare}{铸币贬值与纸币超发}{铸币贬值}{纸币超发}
\cmprow{国家做什么}{减少每枚钱币的含银，或提高钱币的作价}{多印纸币，或多发中央银行的负债}
\cmprow{收入从哪里来}{铸币税 $1-MP/ME$（式~\eqref{eq:medieval-seigniorage}）}{新发的货币能买到的东西（\chref{ch:fiscaldominance}）}
\cmprow{有没有底}{减少含银有：降到零就降不下去了，罗马在 260 年代已接近这个底；提高作价没有}{没有：面值可以一直印大，直到恶性通胀（\chref{ch:hyperinflation}）}
\cmprow{持币的人能做什么}{把含银多的旧钱留下或熔掉，只花新钱；拿钱换实物}{换成外币或实物，尽快把钱花掉}
\cmprow{本书的例子}{罗马三世纪；中世纪欧洲（\chref{ch:medieval}）；德川日本（\chref{ch:tokugawa}）}{宋元明的纸币（\chref{ch:songming}）；魏玛德国（\chref{ch:weimar}）}
\cmpsum{两者都是用发行货币来融资；减少含银有一个金属的底，提高作价与印纸币都没有，所以它们的约束只能来自制度。}
\end{compare}

\section{银少了，物价为什么没有马上涨}\label{sec:greek-prices}

钱里的银子少了，埃及的物价却没有马上涨；274 年钱变好了，物价反而暴涨。一枚钱的价值，到底由什么决定？

\subsubsection{埃及的小麦价格}
先看数据（图~\ref{fig:greek_silver}）。160 年以前，私人小麦价格有 11 份文书，最早的在 45 年。价格在每阿塔巴 5 到 12 德拉克马之间，中位数 8.73。年代正好记为 160 年的还有 2 份。一份是 11，另一份读法有疑问，附录取作 5.52。把这两份算进去，中位数仍是 8.73。此后到 190 年之间没有私人价格。

190—270 年有 27 份文书。中位数是 16 德拉克马，其中 9 份正好是 16。最低 9，最高 28。这 27 份里，26 份出自庄园（一份出自神庙）的账目与书信，20 份出自阿尔西诺伊特州（Arsinoite）\cite{kelly2021}。249—270 年有 18 份。它们的中位数也是 16，而且全都是 4 的倍数。Kelly 说，价格往往以四德拉克马为一档，如 8、12 或 16 德拉克马\cite{kelly2021}，也就是以一枚四德拉克马为单位。190 年以后最早的一份是 192 年的一本庄园账，三个价格的中位数是 18。

这正是 Rathbone 的看法。据 Howgego 等的转述，Rathbone 认为 160—190 年的翻倍是唯一显著的变化。他说的是 45 年到 274/5 年之间的整个时期\cite{howgego2010}。按中位数，190—270 年的价格是 160 年以前的 $16/8.73=1.83$ 倍。

然后是 274 年。Kelly 引 Rathbone（1997）\cite{rathbone1997}的说法：“约 274 年，价格水平突然跳升了十倍或更多”\cite{kelly2021}。按文书，270 年一本庄园账里是 24 德拉克马。276 年官府的折价是 200。这是 190—270 年中位数的 12.5 倍。此后，301 年一张私人收据是 640。314 年一本私人账里是 6000 与 8000。约 340 年，商人申报的是 27 万\cite{kelly2021}。

小麦的价格不是物价水平。小麦有它自己的收成，国家在小麦市场上的分量也比在别的市场上大得多（\secref{sec:greek-test}）。所以这里先把小麦当作一个窗口。物价水平要用 Kelly 合成的生活费用来看，留到\secref{sec:greek-test}。

把价格与钱币放在一起，有三件事要解释。第一，167 年以后，四德拉克马的含银降到原来的 $1/3.83$，小麦价格却只涨到 1.83 倍。按 192 年的价格算，也只是 2.06 倍。第二，190—270 年，罗马银币的含银从四成多降到不足 5\%，埃及的小麦价格却没有趋势性的上涨（\secref{sec:greek-test} 的回归）。埃及用的是自己的四德拉克马，它在这段时间的含银，本章的史料没有逐年的数字。第三，274 年钱变好了，物价反而暴涨。

\fig{greek_silver}{钱币的含银一路下降，小麦价格却先平稳、后暴涨。(a) 看两条线怎样往下走：蓝色是罗马的银币，尼禄 80\%，155 年起 70\%，180 年起 67\%，193 年塞维鲁的标准 44\%，伽利埃努斯时 1\%—2\%（竖条），274 年奥勒良提高到约 5\%（箭头）；193 年与伽利埃努斯之间的点线只表示中间没有逐年的数据。琥珀色虚线是埃及的四德拉克马，18\% $\to$ 5\% $\to$ 10\% $\to$ 5\%。对数纵轴。(b) 看价格在 274 年前后的分别：埃及的小麦价格，每阿塔巴多少德拉克马，对数纵轴，一份文书一个点。朱红色圆点是 274 年以前的私人价格与 274 年以后的全部价格（其中有官方的折价与强制收购），琥珀色方块是 274 年以前的国家价格；两条横线是私人价格的中位数，160 年以前 8.73，190—270 年 16；点线标出 274 年的改革与 301 年的限价令。数据：Howgego 等（2010）表 1—3，Cope（1974），Haklai-Rotenberg（2011）；Kelly（2021）附录 2（CC BY 4.0，2026 年 10 月 14 日下载），一份文书里有几个价格时取中位数，年代取中点。}{fig:greek_silver}

\subsubsection{一枚钱值多少：三种价值}
Haklai-Rotenberg 区分了三种价值，即金属价值（VM）、法定价值（VO）与交换价值（VE）。交换价值不可能低于金属价值，否则持有钱币的人宁可把钱币当作金属块，按金属价值来交易。Haklai-Rotenberg 的结论是，“交换价值从来不低于金属价值，也不高于法定价值”；而且在奥勒良之前的两个世纪里，法定价值的表现是恒定的\cite{haklai2011}。

把 Haklai-Rotenberg 的说法写成一个教学模型，各种价值都折成白银的克数。设一记账单位的钱（例如一德拉克马）含银 $c$ 克，这是它的金属价值；它的法定价值相当于 $\bar v$ 克银，在交易中实际值 $v$ 克银。一阿塔巴小麦值 $q$ 克银，$q$ 变了，是小麦本身变贵或变便宜了。以钱计的小麦价格 $P$，就是 $q$ 除以一德拉克马值多少银。交换价值 $v$ 落在金属价值 $c$ 与法定价值 $\bar v$ 之间。它落在哪里，用 $\tau$ 表示，$\tau$ 衡量的是人们对钱币的信任：
\begin{keyeq}
P=\frac{q}{v},\qquad v=c+\tau\,(\bar v-c),\quad 0\le\tau\le1;\qquad
\eta\equiv-\frac{d\ln P}{d\ln c}=\frac{(1-\tau)\,c}{v} .\label{eq:greek-value}
\end{keyeq}
第一个式子是定义。小麦值 $q$ 克银，一德拉克马值 $v$ 克银，一阿塔巴小麦就要 $q/v$ 德拉克马。第二个式子说，人们愿意为钱币付出的银子，比它所含的多出法定价值与金属价值之差的一个比例 $\tau$。$\tau$ 有两个极端。$\tau=0$ 时，钱只按它所含的银作价，$v=c$，物价 $P=q/c$，这是\chref{ch:tokugawa}的式~\eqref{eq:toku-content}。$\tau=1$ 时，钱按法定价值流通，$P=q/\bar v$，含银多少与物价无关。

第三个式子是\textbf{传导率}，即含银减少 1\%，物价涨百分之几。推导如下。$\bar v$ 与 $\tau$ 不变时，$dv=(1-\tau)\,dc$，两边除以 $v$，得 $d\ln v=(1-\tau)(c/v)\,d\ln c$；$q$ 不变时，$d\ln P=-d\ln v$。$\tau=0$ 时 $\eta=1$，物价与含银一比一反向变动；$\tau=1$ 时 $\eta=0$。在两者之间，钱越是比它的银子值钱（$v/c$ 越大），含银再减少，物价的反应越小。

这个模型有一个要紧的局限。$q$ 与 $v$ 都看不到，看到的只是两者之比 $P$。同样的物价上涨，可以说成钱贬了值，也可以说成小麦变贵了。

\begin{example}{四德拉克马变薄以后}\label{ex:greek-tetra}
按 Howgego 等的表 3，四德拉克马每枚含银 167 年以前 2.3 克，185/6 年以后 0.6 克\cite{howgego2010}。私人小麦价格 160 年以前的中位数是 8.73 德拉克马，192 年是 18 德拉克马。（1）每德拉克马的含银前后各是多少？按成色论（$\tau=0$），192 年的小麦价格应当是多少？（2）实际的传导率（对数变化之比）是多少？（3）设 $q$ 不变，$\bar v$ 等于 167 年以前每德拉克马的含银。192 年一德拉克马实际值多少克银，是它含银的几倍？$\tau$ 是多少？这时含银再减少 1\%，按式~\eqref{eq:greek-value} 物价涨多少？（4）如果 160 年以前到 192 年的上涨全是因为小麦本身变贵了（$q$ 涨到 2.06 倍），$\tau$ 是多少？
\solution （1）167 年以前 $2.3/4=0.575$ 克，185/6 年以后 $0.6/4=0.15$ 克。小麦值 $q=8.73\times0.575=5.02$ 克银；按成色论，192 年的价格应当是 $5.02/0.15=33.5$ 德拉克马，是原来的 3.83 倍。实际是 18，只是原来的 2.06 倍。（2）$\eta=\ln2.06/\ln3.83=0.54$：含银降多少，物价只涨了一半多。（3）$v=5.02/18=0.279$ 克，是含银的 1.86 倍；$\tau=(0.279-0.15)/(0.575-0.15)=0.30$。这时 $\eta=(1-0.30)\times0.15/0.279=0.37$，含银再减 1\%，物价涨 0.37\%。用 190—270 年的中位数 16 代替 18：上涨 1.83 倍，传导率 0.45，$\tau=0.39$。（4）$q$ 涨到 2.06 倍，价格也涨到 2.06 倍，$v$ 没有变，仍是 0.575 克，$\tau=1$。同样的数字，可以说成钱失去了大部分信任（$\tau=0.30$），也可以说成钱的信任丝毫无损、是小麦变贵了（$\tau=1$）。价格数据本身分不开这两种解释。
\end{example}

\subsubsection{160—190 年：瘟疫还是钱}
例~\ref{ex:greek-tetra} 的两种极端，正好是关于二世纪末物价上涨的两种解释。一种是瘟疫。据 Howgego 等，Rathbone 与 Temin 把安东尼瘟疫认定为造成价格结构性变化的外来冲击。Rathbone（2007）说，瘟疫在 166/7 年传入埃及。没有人怀疑，此后十年里它造成了人口总量的大幅损失，大概在 20\%—30\% 左右\cite{howgego2010}。Haklai-Rotenberg 提到，Scheidel（2002）也把二世纪末的价格上涨解释为安东尼瘟疫的长期后果\cite{haklai2011}。

这一说最有力的证据是时间，瘟疫与价格上涨几乎同时。而且按\chref{ch:medieval}的马尔萨斯模型（\secref{sec:medieval-malthus}），人口骤减会抬高工资；按\chref{ch:silver}的数量方程（式~\eqref{eq:silver-qty}），产出 $Y$ 减少而钱 $M$ 不变，物价也要上涨。Kelly 也认为，安东尼瘟疫很可能就是这一变化的导火索\cite{kelly2021}。

另一种是钱。Haklai-Rotenberg 提到，Carrié（2007）认为 180 年代的价格上涨“完全出于货币原因”\cite{haklai2011}。这一说最有力的证据是四德拉克马本身。按 Howgego 等的说法，它的大贬值在 164/5 年就开始了，那时瘟疫还没有到来。含银降到原来的 $1/3.83$，足以解释物价的翻倍还有余。

两说都有弱点。瘟疫说解释了物价的一次跳升，却难以解释 Kelly 发现的此后一百年里其他商品价格的持续上涨。Kelly 说，“一百多年里持续的趋势，说明背后有内生的、系统性的经济驱动力”\cite{kelly2021}。货币说解释了方向，却要回答为什么物价只涨了含银降幅的一半多。Howgego 等的结论是：“无论怎样评价瘟疫，都没有哪一种解释足以说明帝国内部正在发生的重大变化。”\cite{howgego2010}按例~\ref{ex:greek-tetra}，两者各占多少，取决于 $q$ 涨了多少，而这正是史料给不出的。图~\ref{fig:greek_trust} 画出了不同的 $\tau$ 下含银与物价的关系。

\fig{greek_trust}{同样少了银子，物价涨多少取决于信任 $\tau$。看三条线怎样分开：横轴是钱币的含银与原来之比（对数，向右减少），纵轴是物价与原来之比（对数）；实线 $\tau=0$ 是按成色作价，含银减半物价翻倍；虚线 $\tau=0.3$；点线 $\tau=1$ 是按法定价值作价，物价不动。两个方块是例~\ref{ex:greek-tetra} 的埃及：192 年的价格是 160 年以前中位数的 2.06 倍（实心），190—270 年的中位数是 1.83 倍（空心），都落在 $\tau=0.3$ 附近。蓝点是例~\ref{ex:greek-trust} 的 274 年：设改革前 $\tau=1$，改革后的物价落在 $\tau=0$ 的线旁。教学模型，式~\eqref{eq:greek-value}；设法定价值等于原来的含银，$q$ 不变。}{fig:greek_trust}

\subsubsection{274 年：钱变好了，物价涨了十倍}
274 年初，奥勒良发行了两个系列的新币。按 Haklai-Rotenberg 的说法，含银从大约 2\% 或 3\% 提高到将近 5\%\cite{haklai2011}。埃及的四德拉克马也在 274/75 年同时改革，化验显示它在 274 年的含银略有提高\cite{haklai2011}。按成色论，物价应当下降。

Haklai-Rotenberg 的解释是信任。早在 260 年，一份奥克西林库斯纸草（P.Oxy. 1411）就为这种信任的下降提供了证据。那是一份布告，说的是“当地的银钱商拒绝接受新铸的钱币”\cite{haklai2011}。不过 Haklai-Rotenberg 也提醒，这份纸草确定说明的只是奥克西林库斯银钱商的行为；“对钱币的信任虽然因为急剧的贬值而在下降，却远没有丧失”\cite{haklai2011}。

274 年的改革为什么让物价暴涨？Haklai-Rotenberg 推测，一个新的、提高了的法定价值，很可能同时适用于新币和旧币。而“一旦干预发生过，随意提高法定价值这件工具用过了，任何价值都可以强加”\cite{haklai2011}。按这个解释，人们不再相信钱币的作价，只按它所含的银子收钱。

\begin{example}{奥勒良的改革}\label{ex:greek-trust}
取 190—270 年私人小麦价格的中位数 16 德拉克马作为改革前的价格，276 年的 200 德拉克马作为改革后的价格。（1）含银从约 2.5\% 提高到 5\%，按成色论，物价应当怎样变？（2）设法定价值 $\bar v=1$（每单位面值），相当于一枚同重、含银 70\% 的钱币所含的银（Howgego 等表 1 中 155—180 年的标准）；改革后含银 5\%，每单位面值的含银 $c$ 是多少？设改革前人们按法定价值收钱（$\tau=1$，这是对改革前的信任最有利的假设），$q$ 不变。改革后 $v$ 与 $\tau$ 是多少？$\tau=0$ 时物价会涨到多少倍？（3）改革前的 $\tau$ 最低可能是多少？（4）以 270 年的 24 德拉克马作为改革前的价格呢？
\solution （1）每单位面值的含银翻了一倍，按式~\eqref{eq:toku-content} 物价应当减半（$\times0.5$）；实际是 $200/16=12.5$ 倍。（2）$c=5/70=0.071$；物价涨到 12.5 倍，$v=1/12.5=0.08$，$\tau=(0.08-0.071)/(1-0.071)=0.009$，几乎是零：改革后的钱只按所含的银作价。$\tau=0$ 时 $v=c$，物价涨到 $1/0.071=14$ 倍，与实际的 12.5 倍相近。（3）改革后的 $v$ 不能低于 $c$，所以改革前的 $v$ 至少是 $12.5\times0.071=0.89$：改革前，含银只有百分之几的钱，至少按法定价值的九成在流通。（4）$200/24=8.3$ 倍，$v=0.12$，$\tau=0.05$，结论不变。这道题有三个前提，都可疑：$q$ 不变；276 年官府的折价与 270 年以前私人的价格可比；法定价值没有变。如果像 Haklai-Rotenberg 推测的那样，改革同时提高了法定价值，每单位面值的含银还要更少，但提高了多少，本章的史料不能确定。
\end{example}

改革为什么会摧毁信任？Haklai-Rotenberg 指出，奥勒良以后，戴克里先一再试图稳定币制。先是 286 年改变金币的标准，然后是 294/96 年的币制改革。最后，301 年又颁布了币制敕令。这几次看来都没有能恢复稳定。问题也许就出在改革本身：“一再动用币制改革这件工具，本身就很可能造成不信任，应验了公众对政府的货币行为还会反复无常的预期”。直到君士坦丁（Constantine），在奥勒良改革一代多人以后，才建立起一种新的稳定的货币秩序，并维持了几个世纪\cite{haklai2011}。

用式~\eqref{eq:greek-value} 的语言说，$\tau$ 是人们对发行者将来怎样对待这枚钱的预期。一次改革改变的不只是 $c$ 与 $\bar v$，还有这份预期；预期一旦变了，提高含银也换不回来。这与\chref{ch:credibility}讲的政策可信性是同一个问题。

\begin{intuition}[大家都收，我就收]
一枚含银只有百分之几的钱，为什么能按法定价值的九成流通？因为我收下它，是为了把它花出去；只要我相信别人也会按这个价收，我按这个价收就不吃亏。这是一种自我实现的信任：大家都收，它就值钱。反过来，一旦大家都开始怀疑，先拒收的人就少吃亏，信任会一下子垮掉，而不是慢慢减少。怀疑可以早就有了而信任照旧：260 年奥克西林库斯的银钱商拒收过新币，此后埃及的物价又平稳了十几年。这与埃及物价的形状相合：几十年不动，然后几年之内涨十倍。
\end{intuition}

\begin{compare}{钱的价值：成色还是信任}{成色论（金属论）}{信任论（名目论）}
\cmprow{一枚钱值多少}{它所含的银：$v=c$，$\tau=0$}{人们相信它值的：最多到法定价值 $\bar v$，$\tau=1$}
\cmprow{含银减半，物价}{翻一倍（$\eta=1$，式~\eqref{eq:toku-content}）}{只要信任还在，不变（$\eta=0$）}
\cmprow{167—192 年的埃及}{预测涨到 3.83 倍，实际只有 2.06 倍}{实际的上涨只有预测的一半多，钱还有相当的信任}
\cmprow{274 年的改革}{含银提高，物价应当减半}{信任崩溃，物价涨到接近按含银作价的水平}
\cmprow{最难解释的事实}{190—270 年含银大减而小麦价格不涨；274 年的暴涨}{长期看物价终究跟着钱的贬值走：约 340 年的小麦价格是 160 年以前的约 3.1 万倍}
\cmpsum{成色论说钱值它的银子，信任论说钱值人们相信它值的；三世纪的埃及先像后者，274 年以后越来越像前者。}
\end{compare}

\subsubsection{四种解释}
罗马的贬值与物价之间是什么关系？Haklai-Rotenberg 概括过研究者的几种看法，加上这位作者本人的信任说，一共四种，每一种都有它最有力的论证。

第一种可以叫作贬值即通胀。Haklai-Rotenberg 说，有一个研究传统“把三世纪的铸币贬值看作通胀的确证，例如 Jones（1974）”\cite{haklai2011}。它最有力的论证在长期。银子最后几乎没有了，物价最后涨了几万倍，方向完全一致。它的弱点在短期，190—270 年的埃及价格不支持它。

第二种只看价格。Haklai-Rotenberg 这样概括 Strobel（1989）的看法：“贬值本身不足以作为通胀的指标，通胀的证据应当到价格的证据里去找”\cite{haklai2011}。图~\ref{fig:greek_silver}(b) 就是它最有力的论证。

第三种看数量。Haklai-Rotenberg 提到，Callu（1969）认为通胀“不是铸币贬值的必然结果，而是它对钱币数量的影响的结果”。Scheidel（2010）也认为，钱币的质量与数量是影响古代货币制度的最重要的因素\cite{haklai2011}。按\chref{ch:silver}的式~\eqref{eq:silver-qty}，贬值让同样的银子铸出更多的钱。物价涨不涨，要看流通中的钱多了多少，而这在古代几乎无法计量。

第四种是 Haklai-Rotenberg 自己的信任说。274 年的反例最有利于它；它的弱点是信任本身看不见，只能从价格反推，有循环论证的危险。四种看法并不互斥，长期看成色与数量，短期看信任。当时的皇帝又把涨价归咎于谁？

\section{戴克里先的限价令}\label{sec:greek-edict}

戴克里先的回答，写在他那道限价令的序言里。限价令规定了什么？它为什么失败了？

\subsubsection{一千多种最高价}
肯特说，301 年，戴克里先不得不颁布一道法令，为几乎所有的物品与劳务规定最高价格。这道法令“至少有 32 张价目表，一千多个项目”。到肯特写作时，已经发现了 35 个大大小小的抄本，全都来自希腊、小亚细亚与埃及，那是戴克里先自己管辖的地区。所以肯特认为，自然的结论是，它只打算用于戴克里先亲自管辖的那一部分\cite{kent1920}。

序言把物价上涨完全归咎于商人的贪婪。肯特说，似乎最让皇帝痛心的，是士兵来买东西时，物价被抬到最高\cite{kent1920}。序言说，商人加价“不是四倍、八倍，而是达到了人类的语言都找不到合适的词来形容的程度”，有时买一件东西，“士兵的赏赐和军饷就都被抢走了”\cite{kent1920}。限价令的财政动机就在这里。军饷是按钱定的，物价涨了，士兵的实际收入就少了。

法令的办法是规定一个最高限，而不是把商品的价格固定下来。罚则极重：“任何人如果胆敢违反这条法令的文字，都要处以死刑”。急于购买、同违法零售的贪婪商人成交的人，也要承受同样的危险。手里有食用或日用的货物，在这道规定以后却想把它们从市场上撤下来的人，也同样有罪，“因为对制造匮乏的人，惩罚应当比对违反法令利用匮乏的人还要重”\cite{kent1920}。

肯特不接受皇帝的说法。他认为，戴克里先对形势的描述，也许和那个基督徒作家的描述一样片面。除了四个宫廷的开销，还有一个原因，肯特有意留到最后才讲，那就是铸币的贬值。戴克里先以前，皇帝们用旧的名称发行镀锡的铜币，银币和金币自然从流通中消失了。戴克里先干脆“发行了一种新的第纳尔，坦白地就是铜的，并不假装是别的什么”，物价随之“向上调整，而且是大大地向上”\cite{kent1920}。限价令的价格以第纳尔计，肯特就把它当作这种铜第纳尔来折算。

肯特的这几句是 1920 年的概括。Haklai-Rotenberg 把戴克里先的币制分成几步来看。286 年，戴克里先改变了金币的标准。294/96 年，他改革了币制，301 年又颁布了币制敕令（\secref{sec:greek-prices}）。好钱从流通中消失，正是“劣币驱逐良币”。\chref{ch:medieval}说，戴克里先的限价令面对的价格上涨与货币的贬值分不开，肯特的这段话说明了这一点。肯特还说，赋税的结果是一个“恶性循环，不只是税越来越难收，商品的价钱也越来越高”\cite{kent1920}。

限价令之前的价格上涨从何而来？Haklai-Rotenberg 在一条注里逐一排除了几种说法。引发 301 年限价令的价格上涨，不大可能延续自对 25 年前奥勒良币制改革的反应，否则限价令应当颁布得更早。它也可能是对同一年的币制改革的直接反应。可是，“考虑到罗马帝国之内运输与通信所需的时间，301 年的两道法令在时间上如此接近，使人难以接受这个解释”。剩下的可能是，301 年的这次上涨，多半是戴克里先自己 294/96 年改革意外的、不想要的后果\cite{haklai2011}。按 274 年算，301 年其实是 27 年以后。

\begin{example}{一天的工钱买多少大麦}\label{ex:greek-wage}
肯特把限价令的价格折成了 1920 年的美元：“一罗马磅精炼黄金定价 5 万第纳尔，所以一枚铜第纳尔值 0.434 美分”；大麦每蒲式耳 0.873 美元，农业工人每天 0.108 美元，石匠或木匠每天 0.217 美元。肯特在注里说，意大利的莫迪乌斯（modius）合 7.95 干夸脱，军用的莫迪乌斯是它的两倍；一罗马磅合 0.875 金衡磅；可惜小麦的价格“在记录中缺失了”\cite{kent1920}。按美制，1 蒲式耳 = 32 干夸脱，1 干夸脱 = 1.101 升，1 金衡磅 = 12 金衡盎司。（1）大麦在限价令里是每军用莫迪乌斯多少第纳尔？（2）农业工人与石匠一天的工钱各是多少第纳尔？农业工人一天的工钱能买多少大麦？（3）肯特用的金价是每金衡盎司多少美元？（4）限价令的阿埃扎尼（Aezani）抄本给出的金价是一磅 72000 第纳尔\cite{haklai2011}。按这个金价，一第纳尔值多少美分？肯特折算的美元价格要打几折？
\solution （1）一军用莫迪乌斯 $2\times7.95=15.9$ 干夸脱，合 17.5 升。$0.873\times15.9/32/0.00434=99.95$，即 100 第纳尔。（2）农业工人 $0.108/0.00434=24.88$，即 25 第纳尔；石匠 $0.217/0.00434=50$ 第纳尔。农业工人一天的工钱买 $25/100=0.25$ 军用莫迪乌斯大麦，即 3.98 干夸脱，约 4.4 升。（3）一罗马磅 $0.875\times12=10.5$ 金衡盎司，$50000\times0.00434/10.5=20.67$ 美元。（4）$72000/50000=1.44$，一第纳尔值 $0.434/1.44=0.301$ 美分，肯特的美元价格都要乘 0.69。用同一道法令的两个抄本，折算出来的美元价格就差三成；这些美元数字说的是限价令里的比价，不是当时的生活水平。肯特在另一条注里也提醒：“这些工人里，许多人的工钱与食物的价格相比显得很低，但应当注意，他们多数由雇主管饭。”\cite{kent1920}一天的工钱不是工人一天的全部所得。
\end{example}

\subsubsection{限价会造成多大的短缺}
限价令把名义价格钉住，而它管不住的钱在继续贬值。用一个线性的供求模型看。设一种商品的需求是 $A-bp$，供给是 $B+dp$，$p$ 是以 301 年的钱计的实际价格，市场出清的价格是 $p^{*}=(A-B)/(b+d)$。限价令把名义价格定在 $\bar p$，此后价格水平按每年 $\pi$ 上涨，$P_t=(1+\pi)^t$（301 年为 1）。限价在第 $t$ 年的实际水平是 $\bar p/P_t$；它低于 $p^{*}$ 时，需求多于供给，短缺是
\begin{keyeq}
X_t=(A-b\,\bar p/P_t)-(B+d\,\bar p/P_t)=(b+d)\left(p^{*}-\frac{\bar p}{P_t}\right),\qquad P_t=(1+\pi)^t .\label{eq:greek-ceiling}
\end{keyeq}
即使限价一开始定在市场价格上（$\bar p=p^{*}$），只要价格水平还在涨，实际的限价每年都在下降，短缺每年都在扩大。卖出的数量由供给决定，是 $B+d\,\bar p/P_t$，也在下降。限价令的本意是让人买得起，结果是能买到的越来越少。

\begin{example}{限价令的四年}\label{ex:greek-ceiling}
设 $A=200$，$b=1$，$B=0$，$d=1$，限价定在 301 年的市场价格上。价格水平按埃及小麦价格 301—314 年的年均涨幅上涨：从 301 年的 640 德拉克马到 314 年的 7000，13 年涨到 10.9 倍。（1）301 年的市场价格与成交量是多少？（2）年均涨幅 $\pi$ 是多少？（3）一年后，实际限价、需求、供给与短缺各是多少？（4）四年后（305 年，戴克里先退位）呢？
\solution （1）$p^{*}=200/2=100$，成交 100；在这一点上需求的价格弹性是 $1\times100/100=1$。（2）$\pi=(7000/640)^{1/13}-1=20.2\%$。（3）实际限价 $100/1.202=83.2$；需求 116.8，供给 83.2，短缺 33.6。（4）价格水平涨到 $1.202^{4}=2.09$ 倍，实际限价 47.9；需求 152.1，供给 47.9，短缺 104.2，比当初的成交量还多。卖出的数量从 100 降到不到一半（图~\ref{fig:greek_ceiling}）。这里的数字是教学假设，说明的是方向与速度：限价越久，短缺越大。
\end{example}

\fig{greek_ceiling}{名义限价不动，物价在涨，短缺一年比一年大。(a) 看水平线怎样往下移：绿色是需求，琥珀色虚线是供给，横轴是数量，纵轴是以 301 年价格水平计的价格；点线是 301 年定在市场价格上的限价 100，实线是四年后它的实际水平 47.9，两条曲线在这一价格上相差 104。(b) 看柱子怎样长高：限价令以后每年的短缺（柱），以及实际卖出的数量（虚线），设价格水平每年涨 20.2\%（埃及小麦价格 301—314 年的年均涨幅）。教学模型，式~\eqref{eq:greek-ceiling}；需求与供给是本书的假设（例~\ref{ex:greek-ceiling}）。数据：Kelly（2021）附录 2。}{fig:greek_ceiling}

\subsubsection{法令与执行}
限价令执行得怎样？本章只有两类证据，一是一个敌视戴克里先的作家，二是埃及的价格。拉克唐修说，戴克里先用种种勒索使一切东西都变得极其昂贵之后，又试图用一道法令来限制它们的价格。“于是，为了最微不足道的东西流了许多血；人们不敢把任何东西拿出来卖，短缺变得比以前更加严重、更加难以忍受。”他接着说，这道法令在毁了许多人之后，最后出于不得已而被废除了\cite{lactantius}。

肯特的解读是，尽管法令对违犯者规定了死刑，商人们并没有遵守限价令规定的价格。他认为，“流了许多血”最自然的意思，是买不到限价货的人聚众捣毁商铺，杀了商人。但也可能是许多商人，哪怕卖的是不值钱的东西，也被处决了；或者是商人们用武力抗拒官员执行限价，直到流血丧命。其他的商人则把货物囤积起来，等着限制取消的那一天，还在进行的买卖都按非法的价格，偷偷进行\cite{kent1920}。

这正是式~\eqref{eq:greek-ceiling} 预言的短缺，后果是供给退出，交易转入地下。法令维持了多久？肯特先说，它最后不得不废除，生效了多久，没有人知道。然后他推测，它“作为法律，没有超过四年多少，如果超过的话”\cite{kent1920}。

埃及的价格给出了另一类证据。埃及在戴克里先的辖区里，限价令的抄本就有出自埃及的。可是那里的小麦价格照涨不误。301 年，一张私人收据上是 640 德拉克马。到 304/5 年，已是 1200。312/3 年，涨到 2000。314 年，一本私人账里是 6000 与 8000。两者的中位数是 7000\cite{kelly2021}。13 年里，价格涨到 10.9 倍，年均 20.2\%。限价令也许在某些地方、某些时候压住过某些价格，但没有压住埃及的小麦价格。

\begin{pitfall}[罚得够重，限价就能管住]
戴克里先的限价令罚得不能再重：违者处死，买的人同罪，囤积的人也同罪，法令还说对囤积的人理应罚得更重。它还是失败了。原因不在执行不严，而在它管的是结果，不是原因。物价上涨的原因是钱（肯特说的贬值；Haklai-Rotenberg 认为更可能是戴克里先 294/96 年币制改革的后果），限价令没有动钱；按式~\eqref{eq:greek-ceiling}，只要价格水平在涨，名义限价在实际上就越来越低，短缺越来越大。重罚改变不了供求，只改变交易的方式：卖的人把货藏起来，买的人到暗处去买，或者像肯特对拉克唐修的解读那样，聚众砸店。\chref{ch:medieval}的《劳工法令》把工资压在瘟疫以前的水平，结果相似；计划经济里长期的价格管制与排队，见\chref{ch:planned}。
\end{pitfall}

\section{古代经济算不算市场经济}\label{sec:greek-debate}

限价令是一个帝国想要规定上千种价格的尝试。在法令管不到的地方，古代的价格是市场形成的，还是规定的？这也是\chref{ch:mesopotamia}留下的问题。

\subsubsection{Finley 与 Temin}
Temin（2001）说，这场争论有时叫原始主义与现代主义之争，有时叫 Finley 之争。按 Temin 的引述，Finley（1973）\cite{finley1973}宣称：“古代社会没有一个由相互依存的市场组成的庞大集合体那样的经济体系。”Temin 说，Finley 暗中借用了波兰尼的研究\cite{temin2001}，也就是\chref{ch:mesopotamia}里那个说巴比伦没有市场的波兰尼。

Temin 的论点是，“罗马帝国早期的经济主要是一个市场经济”。他也承认，帝国相隔遥远的各部分，并不像今天的市场常常那样紧密地连在一起\cite{temin2001}。他的证据有几类。利率几乎总在 4\% 到 12\% 之间。价格方面，他引 Rathbone 说，小麦、葡萄酒与驴的价格基本上是由自由市场的力量形成的。土地的价格不是固定的。此外还有一次信贷危机。他的结论毫不含糊：“Finley 恰恰错了：古罗马有一个由相互依存的市场组成的庞大集合体那样的经济体系。”\cite{temin2001}

\begin{case}{罗马 · 公元 33 年 · 一场流动性危机}
塔西佗（Tacitus）的《编年史》第六卷记下了这件事。告发者一拥而上，控告那些靠放贷致富的人，说他们违反了独裁官凯撒的一条法律。那条法律关乎借贷的条件与在意大利持有抵押，早已废弛\cite{tacitus}。罗马很早就限制过利息，塔西佗回顾说，十二铜表法规定，“任何人收取的利息不得高于百分之十二”。元老们自己没有一个是清白的，于是向皇帝求得了宽免，债务人与债权人得到一年零六个月来结清所有的账。结果，“钱变得极其稀缺：一方面，所有的债一下子都被催收；另一方面，被定罪的人那么多，他们的财产一拍卖，流通的钱币就被吞进了国库或皇帝的金库”\cite{tacitus}。

元老院又规定，每个债权人要把债款的三分之二投到意大利的土地上。可是放贷的人把所有的钱都囤起来准备买地，要卖的地产又多得很，地价被压得惨不忍睹\cite{tacitus}。最后，皇帝提比略（Tiberius）出手，“借出十万大塞斯特斯，期限三年，不收利息，条件是每个借款人要以价值加倍的地产向人民作抵押”。塞斯特斯是罗马的一种货币单位。“信用就这样恢复了；私人放贷的人也渐渐出现了。”\cite{tacitus}

Temin 把它读作一场流动性危机。利率上升，贷款被催收，地价崩溃；提比略拿出一大笔钱，无息借给地主三年，以恢复流动性\cite{temin2001}。在恐慌中向银行或市场放贷的机构，叫作最后贷款人，通常是国家或中央银行；中央银行怎样做最后贷款人，是\chref{ch:goldstandard}的主题。提比略的贷款有抵押（价值加倍的地产），不收利息，数额事先定好。这场危机说明，罗马已经有了一个信贷市场。钱从流通中被抽走，利率与地价就会动；当局注入流动性，信用就会恢复。它也说明了另一面。危机是由执行一条法律引发的，又是由皇帝的一笔钱化解的，国家在这个市场里不是旁观者。
\end{case}

\subsubsection{Finley 一方最有力的证据}
下面四条不是 Finley 本人的原话，而是从本章的史料里整理出的、支持他这一方的论据。这样的论据也不少，而且多半来自国家。

第一，Kelly 发现，国家的小麦价格“看起来是按惯例定的，通常是 8 德拉克马，而不是通过讨价还价与市场力量形成的”\cite{kelly2021}。274 年以前有 31 个国家价格。其中 20 个正好是 8。最晚的一个 8，出现在 216 年。

第二，国家在小麦上有特别的分量。国有土地上的赋税或租金，通常按小麦计算，大多也用小麦交纳\cite{kelly2021}。在 Kelly 的样本里，小麦价格有 44\% 是国家价格，大麦只有 16\%。缺粮的时候，国家可以动用强制。246 年，“长官下令申报奥克西林库斯州所有的私人存麦”，并可按每阿塔巴 24 德拉克马为国家收购\cite{kelly2021}。

第三，私人的小麦价格也与别的商品不同。Kelly 认为，单价不高的商品，价格带奥波尔的零头，是自由市场交易的迹象。葡萄酒的交易里明显有这种带零头的定价，许多价格带有奥波尔。大麦也有。按附录，190—270 年共有 10 份大麦价格。其中 4 份带零头。小麦却不同，“即使在私人买卖里，带奥波尔零头的定价也比较少见，160 年以后没有一例”\cite{kelly2021}。至于以四德拉克马为一档，Kelly 解释说，银德拉克马通常用亚历山大的四德拉克马币来付，一枚值四德拉克马，所以交易金额往往是四的倍数\cite{kelly2021}。

第四，戴克里先的限价令本身，就是一个帝国试图规定上千种价格的证据。

两边的证据并不矛盾。Kelly 的结论是，埃及有“两个并存的市场部门：一个是小麦，它的价格似乎受到了国家行为的影响；另一个是其他商品，它们在一个相对自由的市场里自己找到市场水平”\cite{kelly2021}。Kelly 也强调，国家对小麦的影响“绝不是像戴克里先限价令那样试图固定价格”\cite{kelly2021}。起作用的是国家作为买者与收税者的分量。国家是经济中的一个主要角色，如果它对涨价反应较慢，或者总在私人价格区间的低端购买，就可能压低平均价格\cite{kelly2021}。埃及的价格里，真能看出这样两个部门吗（\secref{sec:greek-test}）？

\begin{compare}{古代经济：原始主义与市场经济}{原始主义（Finley）}{市场经济（Temin）}
\cmprow{核心论断}{没有“由相互依存的市场组成的庞大集合体”}{“主要是一个市场经济”，只是各地连得不如今天紧}
\cmprow{价格}{由惯例、国家与身份决定：国家的小麦价格多半是 8 德拉克马}{由供求决定：私人价格随时间涨落，利率在 4\%—12\% 之间}
\cmprow{国家}{国家与军队是经济的中心：赋税用小麦交，缺粮时强制征购}{国家是市场里的一个大买者；33 年的危机由市场传导}
\cmprow{最有力的证据}{国家价格的惯例；私人小麦价格 160 年以后没有奥波尔的零头；限价令}{33 年的信贷危机；私人价格与地价的变动}
\cmprow{弱点}{私人价格确实在动，信贷危机也是市场式的}{证据多来自埃及与意大利，国家的分量难以量化；利率的上限 12\% 与塔西佗说的十二铜表法的上限相同}
\cmpsum{两种解读描述的是同一个经济的两个部门：国家深度介入的小麦，和价格自己找水平的其他商品。}
\end{compare}

\section{数据检验与反例}\label{sec:greek-test}

Kelly 的价格可以检验前面的三个说法。190—270 年，小麦价格真的平稳吗？其他商品是不是在涨？从长期看，物价是否终究跟着钱走？

\subsubsection{190—270 年：小麦平稳，大麦在涨}
把私人小麦价格的对数对年份做最小二乘回归，一份文书一个观测（图~\ref{fig:greek_sectors}）。190—270 年有 27 份文书。斜率是每年 $-0.09\%$，95\% 置信区间是 $-0.67\%$ 到 $0.49\%$，$p=0.75$。$R^{2}=0.004$，不取对数时是 0.016。这段时间里看不出趋势，与 Rathbone 的看法一致。Kelly 报告的 $R^{2}$ 是 0.01（$n=27$）\cite{kelly2021}。

大麦不同。同一时期有 10 份大麦价格，斜率是每年 $+1.46\%$，95\% 置信区间是 $0.28\%$ 到 $2.64\%$。$p=0.021$，$R^{2}=0.50$。按这个速度，80 年涨到 3.2 倍。Kelly 报告的 $R^{2}$ 是 0.61（$n=10$）\cite{kelly2021}。我们的数字与 Kelly 的不完全一样，可能是因为年代或中位数的处理不同，也可能是因为拟合的形式不同。Kelly 在注 13 里说，图上的 $R^{2}$ 用的是三次多项式。注 14 的这两个数出自 190 年以后取对数的趋势检验，Kelly 没有说用的是直线还是多项式\cite{kelly2021}。这 10 份里有一份是国家的强制收购（209 年，每阿塔巴 5 德拉克马）。10 个观测本来就少，结论要谨慎。

Kelly 把各种商品合在一起看（表 1）。先看 190—270 年。小麦的年均涨幅是 0.1\%，大麦 1.4\%，葡萄酒 0.7\%。驴是 1.5\%，合成的生活费用是 1.1\%。更早的两段，合成的生活费用涨得慢一些。50—160 年每年涨 0.3\%。160—190 年是 0.8\%。这一段，小麦每年涨 0.7\%\cite{kelly2021}。按 1.1\% 计，80 年涨到 2.40 倍，正是 Kelly 说的“在这段时期里，生活费用会增加到一倍以上”\cite{kelly2021}。Kelly 给出了一条经验法则。罗马埃及的生活费用，可以假定在公元 50 年到 150 年稳定。此后到公元 270 年，平均每年上涨约 1\%\cite{kelly2021}。

\fig{greek_sectors}{两个部门：小麦平稳，大麦在涨，国家的价格多半是定好的。(a) 看两条拟合线的坡度：190—270 年的私人小麦价格（朱红色圆点，27 份文书）与大麦价格（琥珀色方块，10 份），对数纵轴；直线是对数价格对年份的最小二乘拟合，小麦每年 $-0.09\%$（$p=0.75$），大麦每年 $+1.46\%$（$p=0.021$）。(b) 看中间那根柱子：274 年以前的 31 个国家小麦价格，低于 8 德拉克马的 6 个，正好是 8 的 20 个，高于 8 的 5 个（其中有 246 年的强制征购价 24）。数据：Kelly（2021）附录 2（CC BY 4.0，2026 年 10 月 14 日下载）；回归用 statsmodels。}{fig:greek_sectors}

这些结果支持 Kelly 说的两个部门。小麦价格平稳，其他商品在涨。它们也修正了\secref{sec:greek-prices} 的第二个问题。190—270 年，物价水平并非不动，只是涨得慢，每年约 1\%，而小麦恰好不涨。小麦为什么不涨？Kelly 的解释是国家的分量；另一种可能是小麦本身的供给，例如收成。本章的史料分不开这两者。

\subsubsection{长期：物价终究跟着钱走}
把时间拉长，结论就不同了。约 340 年，小麦价格是 27 万德拉克马。这是 160 年以前中位数的约 3.1 万倍。276 年到约 340 年，共 64 年。其间价格从 200 涨到 27 万，年均涨 11.9\%。301 年到 314 年，年均涨 20.2\%。银子没有了，钱的价值也就没有了着落。从几十年、上百年看，物价确实跟着钱的贬值走。不过 274 年改革以后的含银，本章的史料没有数字，所以不能检验物价的涨幅与含银的降幅是否相称。能说的是方向。成色论在短期不灵，在长期方向是对的。

\subsubsection{反例}
有三个事实与本章的某个说法不合。第一，274 年钱变好了，物价反而暴涨（例~\ref{ex:greek-trust}），这是成色论的反例。第二，190—270 年罗马银币的含银大减，埃及的小麦价格却不涨，这是“贬值即通胀”的反例。不过合成的生活费用每年仍涨约 1\%，这个反例只对小麦成立。第三，160—190 年物价的翻倍，可以完全不靠货币来解释（例~\ref{ex:greek-tetra} 的第 4 问）。这提醒我们，信任说里的 $\tau$ 是从价格里反推出来的，而反推要假定 $q$ 不变。

\begin{deeper}[这些数字有多可靠]
\textbf{一份文书一个数。}Kelly 把一份文书里的几个价格合成一个中位数，把年代不确切的文书放在年代范围的中点\cite{kelly2021}。中点可能差出几十年：190—270 年的回归里，有一份文书的年代范围是 202—231 年。年代的误差会让斜率偏向零。

\textbf{正文与附录不一致。}Kelly 的正文说，“公元 45 年 P.Mich. II 123 的数据当作一个点处理，中位数是每阿塔巴 3.0 德拉克马”\cite{kelly2021}；可是附录 2 里，45 年那一行是 P.Mich. II 127（6 个价格，中位数 7.445），P.Mich. II 123 在 46 年，价格 8.73。本章按附录计算；把 45 年那一行换成正文的 3.0，160 年以前的中位数仍是 8.73。

\textbf{少数档案。}190—270 年的 27 份私人小麦价格，26 份出自庄园与神庙的账目与书信，20 份出自阿尔西诺伊特州。这是少数几个大庄园的价格，不一定代表整个埃及，更不一定代表整个帝国；而且 Howgego 等说，埃及是“一种封闭的货币制度”\cite{howgego2010}，它的钱与帝国的钱不同。

\textbf{价格的性质。}274 年以后的价格混杂了私人收据、官府的折价、强制收购与商人的申报：276 年的 200 德拉克马是官府把实物税折成钱的价格，327 年的 17333 是强制收购，约 340 年的 27 万是商人的申报。把 270 年庄园账里的 24 与 276 年官府的 200 相比，比的不是同一种价格。

\textbf{化验的方法。}各家对含银的估计不同，部分是因为方法不同。Haklai-Rotenberg 提醒，Walker 等人的估计“主要依据 X 射线分析，要谨慎对待”；“Cope 所做的破坏性金属分析显示了相似的格局”，最严重的贬值出现在伽利埃努斯独掌帝位的最后几年，“含银降到 2\%—3\%”\cite{haklai2011}。Haklai-Rotenberg 依据的是 Cope（1972）与 Cope、King 等（1997）的结果；Cope 1974 年论文里的标准是每磅 3 斯克鲁普尔与 6 斯克鲁普尔，折合 1.04\% 与 2.08\%\cite{cope1974}。表~\ref{tab:greek-silver} 的 44.4\% 与“约 50\%”、1\%—2\% 与“不到 5\%”，反映的就是这些差别。

\textbf{同一道法令的不同抄本。}肯特按一罗马磅黄金 5 万第纳尔折算美元，阿埃扎尼抄本给的是 7.2 万（例~\ref{ex:greek-wage}）。Haklai-Rotenberg 说引发限价令的价格上涨不大可能延续自“25 年前”的奥勒良改革，按 274 年算是 27 年。这些出入不改变本章的结论，但说明这些数字要连同出处一起读。

\textbf{样本很小。}大麦的回归只有 10 个观测。去掉 209 年的强制收购，9 个观测的斜率是每年 1.35\%（95\% 置信区间 0.19\% 到 2.52\%，$p=0.029$），方向不变，区间仍然很宽。本章的检验能说明方向，不能给出精确的数。

\textbf{单位、纪年与译本。}本章的钱按 1 塔兰特 = 60 弥那 = 6000 德拉克马换算。1 德拉克马 = 6 奥波尔是雅典的比例，也是 Kelly 对埃及早期的折算。阿塔巴是埃及的谷物容量单位，它合多少升，本章的史料没有可以核对的换算，所以不作换算。年代凡未写明公元前的，都在公元以后。古代作者中只有色诺芬注了生卒年，据所用电子书卷首的说明；其他人不注。希腊、拉丁作者的话据英译本转译，是译文的译文。

\textbf{没有取得的。}Finley 的原书没有取得，他的论断据 Temin 的引述；支持他这一方的四条论据，是从本章的史料里整理出来的。
\end{deeper}

\section{今天还剩什么}\label{sec:greek-legacy}

从吕底亚的钱币到戴克里先的限价令，这一章的几条线索逐条看：
\begin{itemize}
\item \textbf{钱的价值来自金属，还是约定。}亚里士多德的两种回答，在本章的模型里成了信任 $\tau$ 的两端。今天的钱没有金属，$\tau$ 就是全部。它取决于发行者将来怎样对待这种钱，这是\chref{ch:credibility}的可信性问题。中世纪的奥雷姆批评君主随意改钱（\chref{ch:medieval}的\secref{sec:medieval-oresme}），问的也是同一个问题，钱是谁的。
\item \textbf{铸币税与财政主导。}罗马的贬值是为了给军队发饷；用发行货币来填补财政缺口，后来以纸币的形式反复出现。中央银行什么时候会被财政拖着走，见\chref{ch:fiscaldominance}；走到尽头的恶性通胀，见\chref{ch:hyperinflation}与\chref{ch:weimar}。
\item \textbf{价格管制。}戴克里先的限价令几乎为所有的物品与劳务规定了最高价（肯特语）。名义限价加上持续的通胀，短缺一年比一年大（式~\eqref{eq:greek-ceiling}）；中世纪的《劳工法令》（\chref{ch:medieval}）、计划经济的价格管制（\chref{ch:planned}）都要面对这个逻辑。
\item \textbf{资源收入怎么花。}劳里昂的白银是分掉还是造船，色诺芬设想靠银矿养活城邦；今天的产油国面对同样的选择，以及资源收入带来的风险（\chref{ch:oil}）。
\item \textbf{国际货币。}雅典的钱出了边境仍是好银子，所以白银本身成了雅典的出口货；哪些钱在国外有人要，见\chref{ch:hierarchy}。
\item \textbf{最后贷款人。}公元 33 年提比略的无息贷款，是一次以加倍的地产为抵押、向市场注入流动性的救助。中央银行在恐慌中怎样放贷，有一条白芝浩原则，见\chref{ch:goldstandard}。
\item \textbf{古代经济是不是市场经济。}埃及的价格说明，市场与国家并存，小麦受国家影响，其他商品的价格自己找水平。古代中国的国家怎样管理粮价与货币，是下一章的问题（\chref{ch:qinhan}）；德川日本的成色与物价（\chref{ch:tokugawa}）、美洲白银与价格革命（\chref{ch:silver}）会再次检验成色论与数量论。
\end{itemize}

\begin{takeaways}
\begin{itemize}
\item 铸币的印记省掉了称量，也让发行者可以少放银子。亚里士多德的书里有两种回答：钱的价值来自金属（《政治学》讲钱的来历），或来自约定（《伦理学》；《政治学》也转述了这种看法）。
\item 雅典把银矿的意外之财用来造船而不是分掉，地米斯托克利用一个依结果而定的借款合同交给富人去办；色诺芬设想国营奴隶出租养活城邦，按他的算术，六千个奴隶一年 60 塔兰特（一年按 360 天）。
\item 罗马银币的含银从 80\% 降到 2\%；贬值的背景是白银减少与军费增加。塞维鲁的一次贬值，按式~\eqref{eq:medieval-seigniorage} 的毛铸币税率是 37\%，前提是交银子的人还没有看穿。
\item 一枚钱的价值介于金属价值与法定价值之间，$v=c+\tau(\bar v-c)$，传导率 $\eta=(1-\tau)c/v$（式~\eqref{eq:greek-value}）。167—192 年埃及四德拉克马的含银降到 $1/3.83$，小麦只涨到 2.06 倍（传导率 0.54）；274 年含银提高，物价却涨到 12.5 倍，接近信任完全丧失（$\tau\approx0$）时的 14 倍。价格本身分不开钱的贬值与小麦的变贵。
\item 限价令罚得极重，却没有动钱；名义限价不变而物价每年涨 20.2\% 时，四年后实际限价只剩 47.9\%，短缺超过当初的成交量（式~\eqref{eq:greek-ceiling}）。埃及的小麦价格 301—314 年涨到 10.9 倍。
\item 190—270 年，埃及的私人小麦价格没有趋势（每年 $-0.09\%$，$p=0.75$），大麦每年涨 1.46\%（$p=0.021$），生活费用每年约 1\%：国家深度介入的小麦与价格自己找水平的其他商品并存，Finley 与 Temin 各说对了一个部门。
\end{itemize}
\end{takeaways}

\begin{selfcheck}
\begin{enumerate}
\item Cope 说，塞维鲁在 193 年把标准降到“每磅银加 15 盎司合金”；伽利埃努斯时降到“每磅 3 斯克鲁普尔”。两种标准各折合含银多少？
\item 按色诺芬的比例（每个奴隶每天净一个奥波尔，一年 360 天），他计划起步时的一千二百个公共奴隶，一年给雅典带来多少塔兰特？
\item 一枚钱的含银减半，设 $\bar v$ 等于原来的含银，$\tau=0.5$，$q$ 不变。按式~\eqref{eq:greek-value}，它的交换价值是 $\bar v$ 的几分之几？物价涨到几倍？传导率（对数变化之比）是多少？
\item 274 年奥勒良提高了钱币的含银，物价却涨了十倍左右。用式~\eqref{eq:greek-value} 说明这怎么可能。成色论能解释吗？
\item 按例~\ref{ex:greek-ceiling} 的需求与供给，如果价格水平每年涨 20\%，限价令颁布两年后，实际限价是多少？短缺是多少？
\item 塞维鲁把第纳尔的含银从 70.6\% 降到 44.4\%。如果铸币厂仍按旧钱的含银收银子，毛铸币税率是多少？
\item Finley 说古代没有“由相互依存的市场组成的庞大集合体”，Temin 说 Finley“恰恰错了”。本章的哪些证据支持前者，哪些支持后者？
\end{enumerate}
\end{selfcheck}

\begin{answers}
\item $12/(12+15)=44.4\%$；$3/288=1.04\%$。
\item $1200\times360/6/6000=12$ 塔兰特。
\item $v=0.5+0.5\times(1-0.5)=0.75$（以 $\bar v$ 为 1）；物价涨到 $1/0.75=1.33$ 倍；传导率 $\ln1.33/\ln2=0.42$，低于成色论的 1。
\item 物价 $P=q/v$，取决于钱的交换价值 $v$，而 $v=c+\tau(\bar v-c)$。改革提高了 $c$，却摧毁了信任：$\tau$ 从接近 1 降到接近 0，$v$ 从接近 $\bar v$ 降到接近 $c$。只要 $c$ 远低于 $\bar v$（改革后每单位面值只有约 7\%），信任的丧失就远远压过含银的提高（例~\ref{ex:greek-trust}）。成色论只看 $c$，预测物价减半，解释不了。
\item 实际限价 $100/1.2^{2}=69.4$；需求 $200-69.4=130.6$，供给 69.4，短缺 61.1。
\item $1-44.4/70.6=37\%$。
\item 支持 Finley：国家的小麦价格多半按惯例定在 8 德拉克马，赋税与租金用小麦交，缺粮时强制征购；私人小麦价格 160 年以后没有奥波尔的零头；戴克里先试图规定上千种价格。支持 Temin：私人价格（大麦、葡萄酒、驴）随时间上涨，利率在 4\%—12\% 之间，地价会变，33 年的危机像一场市场的流动性危机，并由注入流动性化解。Kelly 的两个部门把两者合在一起。
\end{answers}

吕底亚人在金属上打下印记，钱从此可以数，不必再称。印记背后是发行者的担保，罗马的皇帝把这份担保一点点用掉了。银子越放越少，最后，一道以死刑撑腰的限价令也管不住物价。古代中国的国家管物价，用的是另一套工具，有买进卖出粮食的常平仓，有盐铁的专卖，还有铸钱的垄断。在那里，一枚钱值多少，是不是法令说了算？
